\documentclass[aspectratio=169]{beamer}
\usepackage[utf8]{inputenc}
\usepackage[russian]{babel}
\usepackage{graphicx}
\usepackage{listings}
\usepackage{xcolor}
\usepackage{tikz}
\usetikzlibrary{shapes,arrows,positioning}

% Выбор темы
\usetheme{Madrid}
\usecolortheme{default}

% Определение цветов
\definecolor{MyPink}{RGB}{194, 19, 182}
\definecolor{MyBlue}{RGB}{0, 102, 204}
\definecolor{CodeBg}{RGB}{245, 245, 245}
\definecolor{MyGreen}{RGB}{34, 139, 34}

% Настройка цветов темы
\setbeamercolor{structure}{fg=MyBlue}
\setbeamercolor{palette primary}{bg=MyBlue,fg=white}
\setbeamercolor{palette secondary}{bg=MyPink,fg=white}
\setbeamercolor{palette tertiary}{bg=MyBlue!80,fg=white}

% Настройка листингов кода
\lstset{
    basicstyle=\ttfamily\footnotesize,
    backgroundcolor=\color{CodeBg},
    frame=single,
    breaklines=true,
    keywordstyle=\color{blue},
    commentstyle=\color{MyGreen},
    stringstyle=\color{red},
    showstringspaces=false,
    language=TeX
}

% Информация о презентации
\title[Создание постеров в LaTeX]{Создание научных постеров\\с использованием пакета tikzposter}
\subtitle{Отчет по изучению возможностей LaTeX}
\author{Студент группы НПМмд02-24 Викторов Егор}
\institute{РУДН}
\date{\today}

% Логотип (раскомментируйте и укажите путь к файлу)
% \logo{\includegraphics[height=1cm]{logo.png}}

\begin{document}

% Титульный слайд
\begin{frame}
    \titlepage
\end{frame}

% Содержание
\begin{frame}{Содержание}
    \tableofcontents
\end{frame}

% Раздел 1: Введение
\section{Введение}

\begin{frame}{Что такое научный постер?}
    \begin{columns}[T]
        \begin{column}{0.5\textwidth}
            \textbf{Научный постер} — это визуальное представление результатов исследования
            
            \vspace{0.5cm}
            
            \textbf{Используется для:}
            \begin{itemize}
                \item Конференций
                \item Симпозиумов
                \item Защиты работ
                \item Научных выставок
            \end{itemize}
        \end{column}
        
        \begin{column}{0.5\textwidth}
            \begin{tikzpicture}
                \draw[fill=MyBlue!20, draw=MyBlue, thick] (0,0) rectangle (4,5);
                \node at (2,4.5) {\textbf{ЗАГОЛОВОК}};
                \draw[fill=white] (0.2,3.5) rectangle (1.8,4);
                \draw[fill=white] (2.2,3.5) rectangle (3.8,4);
                \draw[fill=white] (0.2,2) rectangle (1.8,3.3);
                \draw[fill=white] (2.2,2) rectangle (3.8,3.3);
                \draw[fill=white] (0.2,0.2) rectangle (3.8,1.8);
                \node[font=\tiny] at (2,0.5) {Выводы};
            \end{tikzpicture}
        \end{column}
    \end{columns}
\end{frame}

\begin{frame}{Цель и задачи работы}
    \begin{block}{Цель}
        Изучение основных возможностей пакета \texttt{tikzposter} и освоение методов создания структурированных научных постеров
    \end{block}
    
    \vspace{0.3cm}
    
    \begin{block}{Задачи}
        \begin{enumerate}
            \item Изучить работу с многоколоночными макетами
            \item Освоить создание и форматирование блоков контента
            \item Научиться вставлять графические элементы
            \item Изучить работу с цветовым оформлением
            \item Освоить добавление заметок и аннотаций
        \end{enumerate}
    \end{block}
\end{frame}

% Раздел 2: Структура постера
\section{Структура постера и колонки}

\begin{frame}[fragile]{Основы работы с колонками}
    \textbf{Окружение columns} позволяет создавать многоколоночные макеты
    
    \vspace{0.3cm}
    
    \begin{lstlisting}
\begin{columns}
  \column{.33}
  % Первая колонка (33% ширины)
  
  \column{.33}
  % Вторая колонка (33% ширины)
  
  \column{.33}
  % Третья колонка (33% ширины)
\end{columns}
    \end{lstlisting}
    
    \vspace{0.3cm}
    
    \begin{alertblock}{Важно!}
        Числовое значение определяет относительную ширину колонки (от 0 до 1)
    \end{alertblock}
\end{frame}

\begin{frame}{Визуализация колонок}
    \begin{center}
        \begin{tikzpicture}[scale=0.8]
            % Три равные колонки
            \draw[fill=MyBlue!30, draw=MyBlue, thick] (0,0) rectangle (3,4);
            \node at (1.5,2) {\textbf{33\%}};
            
            \draw[fill=MyPink!30, draw=MyPink, thick] (3.2,0) rectangle (6.2,4);
            \node at (4.7,2) {\textbf{33\%}};
            
            \draw[fill=MyGreen!30, draw=MyGreen, thick] (6.4,0) rectangle (9.4,4);
            \node at (7.9,2) {\textbf{33\%}};
            
            \node at (4.7,-0.7) {Три колонки равной ширины};
            
            % Две неравные колонки
            \draw[fill=MyBlue!30, draw=MyBlue, thick] (0,-5.5) rectangle (6.5,-1.5);
            \node at (3.25,-3.5) {\textbf{70\%}};
            
            \draw[fill=MyPink!30, draw=MyPink, thick] (6.7,-5.5) rectangle (9.4,-1.5);
            \node at (8.05,-3.5) {\textbf{30\%}};
            
            \node at (4.7,-6.2) {Две колонки разной ширины};
        \end{tikzpicture}
    \end{center}
\end{frame}

\begin{frame}[fragile]{Колонки различной ширины}
    \textbf{Асимметричный макет} для разных типов контента:
    
    \vspace{0.3cm}
    
    \begin{lstlisting}
\begin{columns}
  \column{.7}
  % Широкая колонка (70% ширины)
  \block{Основные результаты}{
    Здесь размещается основной контент,
    графики, таблицы
  }
  
  \column{.3}
  % Узкая колонка (30% ширины)
  \block{Дополнительно}{
    Ссылки, определения, примечания
  }
\end{columns}
    \end{lstlisting}
\end{frame}

\begin{frame}{Практические рекомендации по колонкам}
    \begin{columns}[T]
        \begin{column}{0.5\textwidth}
            \begin{block}{Рекомендации}
                \begin{itemize}
                    \item Сумма ширин ≈ 1.0
                    \item Оптимально: 2-3 колонки
                    \item Широкие — для основного текста
                    \item Узкие — для списков и ссылок
                \end{itemize}
            \end{block}
        \end{column}
        
        \begin{column}{0.5\textwidth}
            \begin{exampleblock}{Примеры пропорций}
                \begin{itemize}
                    \item 0.33 + 0.33 + 0.33 = 0.99 ✓
                    \item 0.5 + 0.5 = 1.0 ✓
                    \item 0.7 + 0.3 = 1.0 ✓
                    \item 0.6 + 0.4 = 1.0 ✓
                \end{itemize}
            \end{exampleblock}
        \end{column}
    \end{columns}
\end{frame}

% Раздел 3: Блоки контента
\section{Работа с блоками контента}

\begin{frame}[fragile]{Базовый синтаксис блоков}
    \textbf{Команда \textbackslash block} создает блок контента:
    
    \vspace{0.3cm}
    
    \begin{lstlisting}
\block{Заголовок блока}{
  Содержимое блока может включать:
  - Текст
  - Списки
  - Формулы
  - Изображения
}
    \end{lstlisting}
    
    \vspace{0.3cm}
    
    \begin{columns}[T]
        \begin{column}{0.5\textwidth}
            \textbf{С заголовком:}
            \begin{lstlisting}
\block{Заголовок}{Текст}
            \end{lstlisting}
        \end{column}
        
        \begin{column}{0.5\textwidth}
            \textbf{Без заголовка:}
            \begin{lstlisting}
\block{}{Текст}
            \end{lstlisting}
        \end{column}
    \end{columns}
\end{frame}

\begin{frame}[fragile]{Добавление заметок к блокам}
    \textbf{Команда \textbackslash note} создает заметку со стрелкой:
    
    \vspace{0.3cm}
    
    \begin{lstlisting}
\block{Заголовок}{
  Основное содержимое блока
  \vspace{2cm}  % Место для заметки
}

\note[targetoffsetx=5cm, 
      targetoffsety=-3cm, 
      width=8cm]{
  Это заметка, указывающая на блок выше
}
    \end{lstlisting}
\end{frame}

\begin{frame}{Параметры заметок}
    \begin{block}{Основные параметры команды \textbackslash note}
        \begin{description}
            \item[\texttt{targetoffsetx}] Горизонтальное смещение стрелки относительно центра блока
            \item[\texttt{targetoffsety}] Вертикальное смещение стрелки
            \item[\texttt{width}] Ширина текстового блока заметки
        \end{description}
    \end{block}
    
    \vspace{0.3cm}
    
    \begin{alertblock}{Не забудьте!}
        Используйте \texttt{\textbackslash vspace} для создания пространства под заметку в основном блоке
    \end{alertblock}
\end{frame}

\begin{frame}{Визуализация заметок}
    \begin{center}
        \begin{tikzpicture}
            % Основной блок
            \draw[fill=MyBlue!20, draw=MyBlue, thick] (0,0) rectangle (8,3);
            \node[anchor=north] at (4,2.8) {\textbf{Основной блок}};
            \node at (4,1.5) {Содержимое блока};
            
            % Заметка
            \draw[fill=MyPink!20, draw=MyPink, thick] (9,1) rectangle (12,2.5);
            \node[font=\small] at (10.5,1.75) {Заметка};
            
            % Стрелка
            \draw[->, thick, MyPink] (9,1.75) -- (8,1.75);
            
            % Подписи параметров
            \draw[<->, dashed] (8,3.5) -- (9,3.5);
            \node[above, font=\tiny] at (8.5,3.5) {targetoffsetx};
            
            \draw[<->, dashed] (8.5,0) -- (8.5,1.75);
            \node[right, font=\tiny] at (8.5,0.875) {targetoffsety};
        \end{tikzpicture}
    \end{center}
\end{frame}

% Раздел 4: Графические элементы
\section{Вставка графических элементов}

\begin{frame}[fragile]{Особенности работы с изображениями}
    \begin{alertblock}{Важное ограничение!}
        В tikzposter \textbf{нельзя} использовать окружение \texttt{figure}
    \end{alertblock}
    
    \vspace{0.3cm}
    
    \textbf{Вместо этого используется:}
    \begin{itemize}
        \item Окружение \texttt{center}
        \item Команда \texttt{\textbackslash captionof\{figure\}\{...\}}
        \item Пакет \texttt{caption} (обязательно!)
    \end{itemize}
    
    \vspace{0.3cm}
    
    \begin{lstlisting}
% В преамбуле:
\usepackage{caption}
    \end{lstlisting}
\end{frame}

\begin{frame}[fragile]{Вставка растровых изображений}
    \begin{lstlisting}
\block{Результаты эксперимента}{
  \begin{center}
    \includegraphics[width=0.8\textwidth]{image.png}
    \captionof{figure}{Подпись к изображению}
  \end{center}
}
    \end{lstlisting}
    
    \vspace{0.3cm}
    
    \begin{block}{Поддерживаемые форматы}
        \begin{itemize}
            \item PNG — для скриншотов и растровой графики
            \item JPG — для фотографий
            \item PDF — для векторной графики
        \end{itemize}
    \end{block}
\end{frame}

\begin{frame}[fragile]{Вставка векторных рисунков TikZ}
    \begin{columns}[T]
        \begin{column}{0.55\textwidth}
            \begin{lstlisting}[basicstyle=\ttfamily\tiny]
\block{Схема}{
  \begin{center}
    \begin{tikzpicture}
      \draw[thick,->] (0,0) -- (4,0) 
        node[right]{$x$};
      \draw[thick,->] (0,0) -- (0,3) 
        node[above]{$y$};
      \draw[blue,thick] (0,0) 
        parabola (3,2);
    \end{tikzpicture}
    \captionof{figure}{График}
  \end{center}
}
            \end{lstlisting}
        \end{column}
        
        \begin{column}{0.45\textwidth}
            \begin{center}
                \begin{tikzpicture}[scale=0.7]
                    \draw[thick,->] (0,0) -- (4,0) node[right]{$x$};
                    \draw[thick,->] (0,0) -- (0,3) node[above]{$y$};
                    \draw[blue,thick] (0,0) parabola (3,2);
                \end{tikzpicture}
            \end{center}
            
            \vspace{0.3cm}
            
            \textbf{Преимущества TikZ:}
            \begin{itemize}
                \item Векторная графика
                \item Масштабируемость
                \item Интеграция с LaTeX
            \end{itemize}
        \end{column}
    \end{columns}
\end{frame}

\begin{frame}{Рекомендации по изображениям}
    \begin{columns}[T]
        \begin{column}{0.5\textwidth}
            \begin{block}{Для печати}
                \begin{itemize}
                    \item Разрешение ≥ 300 DPI
                    \item Векторные форматы (PDF, EPS)
                    \item Крупные шрифты на графиках
                    \item Высокий контраст
                \end{itemize}
            \end{block}
        \end{column}
        
        \begin{column}{0.5\textwidth}
            \begin{block}{Для экрана}
                \begin{itemize}
                    \item Разрешение ≥ 150 DPI
                    \item PNG с прозрачностью
                    \item Оптимизация размера
                    \item Яркие цвета
                \end{itemize}
            \end{block}
        \end{column}
    \end{columns}
    
    \vspace{0.3cm}
    
    \begin{alertblock}{Совет}
        Всегда проверяйте, как выглядят изображения в финальном размере постера!
    \end{alertblock}
\end{frame}

% Раздел 5: Цветовое оформление
\section{Работа с цветовым оформлением}

\begin{frame}[fragile]{Ограничения tikzposter}
    \begin{alertblock}{Важное ограничение!}
        Пакет tikzposter \textbf{не поддерживает} опцию \texttt{svgnames} пакета xcolor
    \end{alertblock}
    
    \vspace{0.3cm}
    
    \textbf{Это означает, что недоступны:}
    \begin{itemize}
        \item \texttt{Red}, \texttt{Blue}, \texttt{Green}
        \item \texttt{ForestGreen}, \texttt{RoyalBlue}
        \item Другие именованные цвета SVG
    \end{itemize}
    
    \vspace{0.3cm}
    
    \begin{block}{Решение}
        Все цвета должны быть определены вручную с помощью \texttt{\textbackslash definecolor}
    \end{block}
\end{frame}

\begin{frame}[fragile]{Определение пользовательских цветов}
    \textbf{В преамбуле документа:}
    
    \vspace{0.3cm}
    
    \begin{lstlisting}
\definecolor{MyPink}{RGB}{194, 19, 182}
\definecolor{MyBlue}{RGB}{0, 102, 204}
\definecolor{MyGreen}{RGB}{34, 139, 34}
\definecolor{MyOrange}{RGB}{255, 140, 0}
    \end{lstlisting}
    
    \vspace{0.3cm}
    
    \begin{block}{Цветовые модели}
        \begin{description}
            \item[RGB] Красный, Зеленый, Синий (0-255)
            \item[HTML] Шестнадцатеричный код (\#RRGGBB)
            \item[cmyk] Голубой, Пурпурный, Желтый, Черный (0-1)
        \end{description}
    \end{block}
\end{frame}

\begin{frame}[fragile]{Использование цветов в тексте}
    \begin{lstlisting}
Обычный черный текст, 
\color{MyPink} розовый текст, 
\color{black} снова черный текст.
    \end{lstlisting}
    
    \vspace{0.5cm}
    
    \textbf{Результат:}
    
    Обычный черный текст, \color{MyPink}розовый текст\color{black}, снова черный текст.
    
    \vspace{0.5cm}
    
    \begin{alertblock}{Важно!}
        После использования \texttt{\textbackslash color} все последующий текст будет этого цвета, пока не будет задан другой цвет
    \end{alertblock}
\end{frame}

\begin{frame}{Демонстрация цветов}
    \begin{center}
        \begin{tikzpicture}
            % Розовый
            \fill[MyPink] (0,0) rectangle (2,1.5);
            \node[white] at (1,0.75) {\textbf{MyPink}};
            \node[below, font=\tiny] at (1,0) {RGB(194,19,182)};
            
            % Синий
            \fill[MyBlue] (2.5,0) rectangle (4.5,1.5);
            \node[white] at (3.5,0.75) {\textbf{MyBlue}};
            \node[below, font=\tiny] at (3.5,0) {RGB(0,102,204)};
            
            % Зеленый
            \fill[MyGreen] (5,0) rectangle (7,1.5);
            \node[white] at (6,0.75) {\textbf{MyGreen}};
            \node[below, font=\tiny] at (6,0) {RGB(34,139,34)};
            
            % Оранжевый
            \fill[orange] (7.5,0) rectangle (9.5,1.5);
            \node[white] at (8.5,0.75) {\textbf{Orange}};
            \node[below, font=\tiny] at (8.5,0) {RGB(255,140,0)};
        \end{tikzpicture}
    \end{center}
    
    \vspace{0.5cm}
    
    \begin{block}{Инструменты для подбора цветов}
        \begin{itemize}
            \item Google Color Picker
            \item Adobe Color (color.adobe.com)
            \item Coolors.co — генератор палитр
        \end{itemize}
    \end{block}
\end{frame}

\begin{frame}{Принципы работы с цветом}
    \begin{columns}[T]
        \begin{column}{0.5\textwidth}
            \begin{block}{Правило 60-30-10}
                \begin{itemize}
                    \item 60\% — основной цвет
                    \item 30\% — дополнительный
                    \item 10\% — акцентный
                \end{itemize}
            \end{block}
            
            \vspace{0.3cm}
            
            \begin{block}{Контрастность}
                \begin{itemize}
                    \item Темный текст на светлом фоне
                    \item Светлый текст на темном фоне
                    \item Избегайте низкого контраста
                \end{itemize}
            \end{block}
        \end{column}
        
        \begin{column}{0.5\textwidth}
            \begin{exampleblock}{Хорошие сочетания}
                \begin{itemize}
                    \item Синий + оранжевый
                    \item Фиолетовый + желтый
                    \item Зеленый + красный
                \end{itemize}
            \end{exampleblock}
            
            \vspace{0.3cm}
            
            \begin{alertblock}{Избегайте}
                \begin{itemize}
                    \item Более 4-5 цветов
                    \item Слишком яркие цвета
                    \item Плохой контраст
                \end{itemize}
            \end{alertblock}
        \end{column}
    \end{columns}
\end{frame}

% Раздел 6: Полный пример
\section{Полный пример постера}

\begin{frame}[fragile]{Структура документа}
    \begin{lstlisting}[basicstyle=\ttfamily\tiny]
\documentclass[25pt, a0paper, portrait]{tikzposter}
\usepackage[utf8]{inputenc}
\usepackage[russian]{babel}
\usepackage{caption}

\definecolor{MyPink}{RGB}{194, 19, 182}
\definecolor{MyBlue}{RGB}{0, 102, 204}

\title{Название исследования}
\author{Автор}
\institute{Университет}

\begin{document}
\maketitle

% Содержимое постера

\end{document}
    \end{lstlisting}
\end{frame}

\begin{frame}[fragile]{Пример содержимого постера (1/2)}
    \begin{lstlisting}[basicstyle=\ttfamily\tiny]
\begin{columns}
  \column{.5}
  \block{Введение}{
    Краткое описание проблемы и актуальности 
    исследования. Постановка задачи.
  }
  
  \block{Методы}{
    \begin{itemize}
      \item Описание используемых методов
      \item Экспериментальная установка
      \item Параметры эксперимента
    \end{itemize}
  }
    \end{lstlisting}
\end{frame}

\begin{frame}[fragile]{Пример содержимого постера (2/2)}
    \begin{lstlisting}[basicstyle=\ttfamily\tiny]
  \column{.5}
  \block{Результаты}{
    \begin{center}
      \includegraphics[width=0.9\textwidth]{results.png}
      \captionof{figure}{Основные результаты эксперимента}
    \end{center}
  }
  
  \block{Выводы}{
    \begin{enumerate}
      \item Первый важный вывод
      \item Второй важный вывод
      \item Перспективы дальнейших исследований
    \end{enumerate}
  }
\end{columns}
    \end{lstlisting}
\end{frame}

% Раздел 7: Практические рекомендации
\section{Практические рекомендации}

\begin{frame}{Планирование структуры}
    \begin{block}{Перед началом работы}
        \begin{enumerate}
            \item Определите основные разделы постера
            \item Спланируйте количество и ширину колонок
            \item Оцените необходимое пространство для графиков
            \item Продумайте цветовую схему
            \item Подготовьте все изображения в нужном формате
        \end{enumerate}
    \end{block}
    
    \vspace{0.3cm}
    
    \begin{alertblock}{Совет}
        Нарисуйте эскиз постера на бумаге перед началом работы в LaTeX!
    \end{alertblock}
\end{frame}

\begin{frame}{Оформление контента}
    \begin{columns}[T]
        \begin{column}{0.5\textwidth}
            \begin{block}{Текст}
                \begin{itemize}
                    \item Краткие заголовки
                    \item Минимум текста
                    \item Крупный шрифт
                    \item Структурированность
                \end{itemize}
            \end{block}
        \end{column}
        
        \begin{column}{0.5\textwidth}
            \begin{block}{Визуализация}
                \begin{itemize}
                    \item Больше графиков
                    \item Меньше таблиц
                    \item Яркие цвета
                    \item Понятные подписи
                \end{itemize}
            \end{block}
        \end{column}
    \end{columns}
    
    \vspace{0.3cm}
    
    \begin{exampleblock}{Правило «3 метра»}
        Основная информация должна читаться с расстояния 3 метров
    \end{exampleblock}
\end{frame}

\begin{frame}{Типичные ошибки}
    \begin{alertblock}{Чего следует избегать}
        \begin{itemize}
            \item Слишком много текста — постер превращается в статью
            \item Мелкий шрифт — информация нечитаема
            \item Слишком много цветов — пестрота и хаос
            \item Низкое качество изображений — непрофессиональный вид
            \item Отсутствие структуры — сложно следить за логикой
            \item Неправильные пропорции колонок — неэффективное использование пространства
        \end{itemize}
    \end{alertblock}
\end{frame}

\begin{frame}{Технические аспекты}
    \begin{block}{Обязательные требования}
        \begin{itemize}
            \item Подключить пакет \texttt{caption} для работы с подписями
            \item Определить все используемые цвета в преамбуле
            \item Проверить сумму ширин колонок (должна быть ≈ 1.0)
            \item Использовать изображения в высоком разрешении
            \item Добавить \texttt{\textbackslash vspace} для заметок
        \end{itemize}
    \end{block}
    
    \vspace{0.3cm}
    
    \begin{exampleblock}{Размеры постеров}
        \begin{itemize}
            \item A0 — 841 × 1189 мм (стандарт для конференций)
            \item A1 — 594 × 841 мм (компактный вариант)
            \item Ориентация: portrait (вертикальная) или landscape (горизонтальная)
        \end{itemize}
    \end{exampleblock}
\end{frame}

% Раздел 8: Заключение
\section{Заключение}

\begin{frame}{Освоенные навыки}
    \begin{columns}[T]
        \begin{column}{0.5\textwidth}
            \textbf{Структурирование:}
            \begin{itemize}
                \item Многоколоночные макеты
                \item Блоки контента
                \item Заметки и аннотации
            \end{itemize}
            
            \vspace{0.3cm}
            
            \textbf{Оформление:}
            \begin{itemize}
                \item Цветовые схемы
                \item Пользовательские цвета
                \item Форматирование текста
            \end{itemize}
        \end{column}
        
        \begin{column}{0.5\textwidth}
            \textbf{Графика:}
            \begin{itemize}
                \item Растровые изображения
                \item Векторная графика TikZ
                \item Подписи к рисункам
            \end{itemize}
            
            \vspace{0.3cm}
            
            \textbf{Техническое:}
            \begin{itemize}
                \item Работа с пакетами
                \item Настройка параметров
                \item Решение проблем
            \end{itemize}
        \end{column}
    \end{columns}
\end{frame}

\begin{frame}{Преимущества tikzposter}
    \begin{block}{Почему стоит использовать tikzposter?}
        \begin{enumerate}
            \item \textbf{Профессиональное оформление} — без графических редакторов
            \item \textbf{Гибкость} — полный контроль над структурой и дизайном
            \item \textbf{Интеграция} — работает со всеми пакетами LaTeX
            \item \textbf{Масштабируемость} — легко изменить размер постера
            \item \textbf{Воспроизводимость} — легко обновлять и модифицировать
            \item \textbf{Качество} — векторная графика для печати
        \end{enumerate}
    \end{block}
\end{frame}

\begin{frame}{Области применения}
    \begin{columns}[T]
        \begin{column}{0.5\textwidth}
            \begin{block}{Научная сфера}
                \begin{itemize}
                    \item Конференции
                    \item Симпозиумы
                    \item Семинары
                    \item Защиты диссертаций
                \end{itemize}
            \end{block}
        \end{column}
        
        \begin{column}{0.5\textwidth}
            \begin{block}{Образование}
                \begin{itemize}
                    \item Курсовые работы
                    \item Дипломные проекты
                    \item Студенческие конференции
                    \item Научные кружки
                \end{itemize}
            \end{block}
        \end{column}
    \end{columns}
    
    \vspace{0.5cm}
    
    \begin{center}
        \Large\textbf{tikzposter — мощный инструмент для\\визуализации научных результатов!}
    \end{center}
\end{frame}

\begin{frame}{Дальнейшее изучение}
    \begin{block}{Рекомендуемые ресурсы}
        \begin{itemize}
            \item Официальная документация tikzposter
            \item Галерея примеров на Overleaf
            \item Сообщество TeX Stack Exchange
            \item Книга «The TikZ and PGF Manual»
        \end{itemize}
    \end{block}
    
    \vspace{0.3cm}
    
    \begin{block}{Дополнительные возможности для изучения}
        \begin{itemize}
            \item Создание собственных тем оформления
            \item Продвинутая работа с TikZ
            \item Интеграция с другими пакетами (pgfplots, circuitikz)
            \item Автоматизация создания постеров
        \end{itemize}
    \end{block}
\end{frame}

% Финальный слайд
\begin{frame}[plain]
    \begin{center}
        \Huge\textbf{Спасибо за внимание!}
        \vspace{1cm}
    \end{center}
\end{frame}

\end{document}
